% ==================== THEME 2 — AFFECTATIONS / EXPRESSIONS ====================

% ---- aff-egal-vs-double-egal ----
\element{bank-aff-egal-vs-double-egal}{
    \begin{question}{aff-egal-vs-double-egal-v1}
        Pour affecter la valeur \texttt{10} à la variable \texttt{x}, quelle écriture est correcte ?
        \begin{multicols}{4}
            \begin{reponses}
                \bonne{\lstinline[language=C]|x = 10;|}
                \mauvaise{\lstinline[language=C]|x == 10;|}
                \mauvaise{\lstinline[language=C]|10 = x;|}
                \mauvaise{\lstinline[language=C]|x := 10;|}
            \end{reponses}
        \end{multicols}
    \end{question}
}
\element{bank-aff-egal-vs-double-egal}{
    \begin{question}{aff-egal-vs-double-egal-v2}
        Pour affecter la valeur \texttt{25} à la variable \texttt{n}, quelle écriture est correcte ?
        \begin{multicols}{4}
            \begin{reponses}
                \bonne{\lstinline[language=C]|n = 25;|}
                \mauvaise{\lstinline[language=C]|n == 25;|}
                \mauvaise{\lstinline[language=C]|25 = n;|}
                \mauvaise{\lstinline[language=C]|n := 25;|}
            \end{reponses}
        \end{multicols}
    \end{question}
}
\element{bank-aff-egal-vs-double-egal}{
    \begin{question}{aff-egal-vs-double-egal-v3}
        Pour affecter la valeur \texttt{7} à la variable \texttt{score}, quelle écriture est correcte ?
        \begin{multicols}{4}
            \begin{reponses}
                \bonne{\lstinline[language=C]|score = 7;|}
                \mauvaise{\lstinline[language=C]|score == 7;|}
                \mauvaise{\lstinline[language=C]|7 = score;|}
                \mauvaise{\lstinline[language=C]|score := 7;|}
            \end{reponses}
        \end{multicols}
    \end{question}
}
\element{affectations}{
    \restituegroupe[1]{bank-aff-egal-vs-double-egal}
}

% ---- aff-point-virgule ----
\element{bank-aff-point-virgule}{
    \begin{question}{aff-point-virgule-v1}
        Quelle instruction est syntaxiquement correcte en C ?
        \begin{multicols}{3}
            \begin{reponses}
                \mauvaise{\lstinline[language=C]|x = x + 1|}
                \bonne{\lstinline[language=C]|x = x + 1;|}
                \mauvaise{\lstinline[language=C]|x = x + 1;;;|}
            \end{reponses}
        \end{multicols}
    \end{question}
}
\element{bank-aff-point-virgule}{
    \begin{question}{aff-point-virgule-v2}
        Quelle instruction est syntaxiquement correcte en C ?
        \begin{multicols}{3}
            \begin{reponses}
                \mauvaise{\lstinline[language=C]|n = n - 2|}
                \bonne{\lstinline[language=C]|n = n - 2;|}
                \mauvaise{\lstinline[language=C]|n = n - 2;;;|}
            \end{reponses}
        \end{multicols}
    \end{question}
}
\element{bank-aff-point-virgule}{
    \begin{question}{aff-point-virgule-v3}
        Quelle instruction est syntaxiquement correcte en C ?
        \begin{multicols}{3}
            \begin{reponses}
                \mauvaise{\lstinline[language=C]|total = total + p|}
                \bonne{\lstinline[language=C]|total = total + p;|}
                \mauvaise{\lstinline[language=C]|total = total + p;;;|}
            \end{reponses}
        \end{multicols}
    \end{question}
}
\element{affectations}{
    \restituegroupe[1]{bank-aff-point-virgule}
}

% ---- aff-incrementation ----
\element{bank-aff-incrementation}{
    \begin{question}{aff-incrementation-v1}
        Après \lstinline[language=C]|int x = 5; x++;|, quelle est la valeur de \texttt{x} ?
        \begin{multicols}{4}
            \begin{reponses}
                \mauvaise{4}
                \bonne{6}
                \mauvaise{5}
                \mauvaise{Erreur de compilation}
            \end{reponses}
        \end{multicols}
    \end{question}
}
\element{bank-aff-incrementation}{
    \begin{question}{aff-incrementation-v2}
        Après \lstinline[language=C]|int x = 9; x++;|, quelle est la valeur de \texttt{x} ?
        \begin{multicols}{4}
            \begin{reponses}
                \mauvaise{8}
                \bonne{10}
                \mauvaise{9}
                \mauvaise{Erreur de compilation}
            \end{reponses}
        \end{multicols}
    \end{question}
}
\element{bank-aff-incrementation}{
    \begin{question}{aff-incrementation-v3}
        Après \lstinline[language=C]|int x = 2; x++;|, quelle est la valeur de \texttt{x} ?
        \begin{multicols}{4}
            \begin{reponses}
                \mauvaise{1}
                \bonne{3}
                \mauvaise{2}
                \mauvaise{Erreur de compilation}
            \end{reponses}
        \end{multicols}
    \end{question}
}
\element{affectations}{
    \restituegroupe[1]{bank-aff-incrementation}
}

% ---- aff-division-entiere ----
% Durci : la division entière positive (arrondi vers le bas) est déjà bien
% connue en fin de séquence. On introduit un opérande négatif : le piège
% classique est de croire que le C arrondit vers le bas (comme une division
% mathématique/floor), alors qu'il tronque vers zéro.
\element{bank-aff-division-entiere}{
    \begin{question}{aff-division-entiere-v1}
        Que vaut \texttt{r} après \lstinline[language=C]|int r = -7 / 2;| ?
        \begin{multicols}{5}
            \begin{reponses}
                \bonne{-3}
                \mauvaise{-4}
                \mauvaise{-3.5}
                \mauvaise{3}
                \mauvaise{-7}
            \end{reponses}
        \end{multicols}
    \end{question}
}
\element{bank-aff-division-entiere}{
    \begin{question}{aff-division-entiere-v2}
        Que vaut \texttt{r} après \lstinline[language=C]|int r = -9 / 2;| ?
        \begin{multicols}{5}
            \begin{reponses}
                \bonne{-4}
                \mauvaise{-5}
                \mauvaise{-4.5}
                \mauvaise{4}
                \mauvaise{-9}
            \end{reponses}
        \end{multicols}
    \end{question}
}
\element{bank-aff-division-entiere}{
    \begin{question}{aff-division-entiere-v3}
        Que vaut \texttt{r} après \lstinline[language=C]|int r = -11 / 4;| ?
        \begin{multicols}{5}
            \begin{reponses}
                \bonne{-2}
                \mauvaise{-3}
                \mauvaise{-2.75}
                \mauvaise{2}
                \mauvaise{-11}
            \end{reponses}
        \end{multicols}
    \end{question}
}
\element{affectations}{
    \restituegroupe[1]{bank-aff-division-entiere}
}

% ---- aff-modulo ----
\element{bank-aff-modulo}{
    \begin{question}{aff-modulo-v1}
        Que vaut \texttt{r} après \texttt{int r = 13 \% 5;} ?
        \begin{multicols}{4}
            \begin{reponses}
                \mauvaise{2.6}
                \mauvaise{2}
                \bonne{3}
                \mauvaise{5}
            \end{reponses}
        \end{multicols}
    \end{question}
}
\element{bank-aff-modulo}{
    \begin{question}{aff-modulo-v2}
        Que vaut \texttt{r} après \texttt{int r = 17 \% 4;} ?
        \begin{multicols}{4}
            \begin{reponses}
                \mauvaise{4.25}
                \mauvaise{4}
                \bonne{1}
                \mauvaise{5}
            \end{reponses}
        \end{multicols}
    \end{question}
}
\element{bank-aff-modulo}{
    \begin{question}{aff-modulo-v3}
        Que vaut \texttt{r} après \texttt{int r = 22 \% 6;} ?
        \begin{multicols}{4}
            \begin{reponses}
                \mauvaise{3.6}
                \mauvaise{3}
                \bonne{4}
                \mauvaise{6}
            \end{reponses}
        \end{multicols}
    \end{question}
}
\element{affectations}{
    \restituegroupe[1]{bank-aff-modulo}
}

% ---- aff-trace (question de trace : tableau structuré au lieu du quadrillage) ----
\element{bank-aff-trace}{
    \begin{question}{aff-trace-v1}
        Pour chacune des lignes suivantes, indiquer la variable modifiée et sa nouvelle valeur.
        \lstinputlisting[language=c,basicstyle=\Large]{sources/codes/affectation-1.c}
        \evaluationProfTableauTrace{2}{%
            \begin{center}
            \begin{tabular}{|c|c|p{4cm}|}
                \hline
                \textbf{Ligne} & \textbf{Variable} & \textbf{Nouvelle valeur} \\
                \hline
                \texttt{val\_i = val\_i + 5;} & & \\[0.3cm]
                \hline
                \texttt{val\_f = val\_i / 2;} & & \\[0.3cm]
                \hline
                \texttt{val\_f = val\_i / 2.0;} & & \\[0.3cm]
                \hline
                \texttt{val\_c = 'C';} & & \\[0.3cm]
                \hline
                \texttt{val\_c++;} & & \\[0.3cm]
                \hline
                \texttt{val\_i = 17 \% 4;} & & \\[0.3cm]
                \hline
                \texttt{val\_i = 20 \% 4;} & & \\[0.3cm]
                \hline
            \end{tabular}
            \end{center}
        }
    \end{question}
}
\element{bank-aff-trace}{
    \begin{question}{aff-trace-v2}
        Pour chacune des lignes suivantes, indiquer la variable modifiée et sa nouvelle valeur.
        \lstinputlisting[language=c,basicstyle=\Large]{sources/codes/affectation-2.c}
        \evaluationProfTableauTrace{2}{%
            \begin{center}
            \begin{tabular}{|c|c|p{4cm}|}
                \hline
                \textbf{Ligne} & \textbf{Variable} & \textbf{Nouvelle valeur} \\
                \hline
                \texttt{val\_i = val\_i + 4;} & & \\[0.3cm]
                \hline
                \texttt{val\_f = val\_i / 2;} & & \\[0.3cm]
                \hline
                \texttt{val\_f = val\_i / 2.0;} & & \\[0.3cm]
                \hline
                \texttt{val\_c = 'D';} & & \\[0.3cm]
                \hline
                \texttt{val\_c++;} & & \\[0.3cm]
                \hline
                \texttt{val\_i = 19 \% 5;} & & \\[0.3cm]
                \hline
                \texttt{val\_i = 22 \% 5;} & & \\[0.3cm]
                \hline
            \end{tabular}
            \end{center}
        }
    \end{question}
}
\element{bank-aff-trace}{
    \begin{question}{aff-trace-v3}
        Pour chacune des lignes suivantes, indiquer la variable modifiée et sa nouvelle valeur.
        \lstinputlisting[language=c,basicstyle=\Large]{sources/codes/affectation-3.c}
        \evaluationProfTableauTrace{2}{%
            \begin{center}
            \begin{tabular}{|c|c|p{4cm}|}
                \hline
                \textbf{Ligne} & \textbf{Variable} & \textbf{Nouvelle valeur} \\
                \hline
                \texttt{val\_i = val\_i + 3;} & & \\[0.3cm]
                \hline
                \texttt{val\_f = val\_i / 2;} & & \\[0.3cm]
                \hline
                \texttt{val\_f = val\_i / 2.0;} & & \\[0.3cm]
                \hline
                \texttt{val\_c = 'E';} & & \\[0.3cm]
                \hline
                \texttt{val\_c++;} & & \\[0.3cm]
                \hline
                \texttt{val\_i = 23 \% 6;} & & \\[0.3cm]
                \hline
                \texttt{val\_i = 25 \% 6;} & & \\[0.3cm]
                \hline
            \end{tabular}
            \end{center}
        }
    \end{question}
}
\element{affectations}{
    \restituegroupe[1]{bank-aff-trace}
}
